\subsection*{A.3. Đánh giá mô hình (Ch.\,3, Bài tập 01/08)}
Confusion matrix: hàng = nhãn thực, cột = nhãn dự đoán.

\kyhieubatdau
  \kh{n}{tổng số mẫu}
  \kh{\mathrm{TP}_k}{true positive lớp \(k\): thực và đoán đều là \(k\)}
  \kh{\mathrm{TP},\mathrm{TN},\mathrm{FP},\mathrm{FN}}{true positive, true negative, false positive, false negative (nhị phân)}
\kyhieuketthuc

\paragraph{Đa lớp.}
\begin{align*}
\textit{accuracy} &= \frac{\sum_k \mathrm{TP}_k}{n},\\
\mathrm{Precision}_k &= \frac{\mathrm{TP}_k}{\text{tổng cột đoán là } k},\qquad
\mathrm{Recall}_k=\frac{\mathrm{TP}_k}{\text{tổng hàng thực là } k},\\
\mathrm{F1}_k &= \frac{2\,\mathrm{Precision}_k\,\mathrm{Recall}_k}{\mathrm{Precision}_k+\mathrm{Recall}_k}.
\end{align*}
\ynghia{Đo chất lượng phân lớp từng lớp và toàn cục: accuracy = tỉ lệ đoán đúng; precision = độ tin cậy khi đoán lớp \(k\); recall = tỉ lệ bắt được mẫu thực lớp \(k\); F1 = trung bình điều hoà precision và recall.}
\kyhieubatdau
  \kh{k}{chỉ số lớp}
  \kh{\mathrm{TP}_k}{số mẫu thực và đoán đều thuộc lớp \(k\)}
  \kh{n}{tổng số mẫu}
  \kh{\mathrm{Precision}_k}{độ chính xác khi dự đoán lớp \(k\)}
  \kh{\mathrm{Recall}_k}{độ bao phủ (sensitivity) lớp \(k\)}
  \kh{\mathrm{F1}_k}{F1-score lớp \(k\)}
\kyhieuketthuc

\paragraph{Nhị phân (lớp dương).}
\begin{align*}
\textit{accuracy} &= \frac{\mathrm{TP}+\mathrm{TN}}{n},\qquad
\textit{precision}=\frac{\mathrm{TP}}{\mathrm{TP}+\mathrm{FP}},\\
\textit{recall}=\mathrm{TPR}&=\frac{\mathrm{TP}}{\mathrm{TP}+\mathrm{FN}},\qquad
\textit{specificity}=\mathrm{TNR}=\frac{\mathrm{TN}}{\mathrm{TN}+\mathrm{FP}},\\
\mathrm{F1}&=\frac{2\,\textit{precision}\,\textit{recall}}{\textit{precision}+\textit{recall}},\qquad
\textit{prevalence}=\frac{\mathrm{TP}+\mathrm{FN}}{n}.
\end{align*}
\ynghia{Các metric chuẩn cho bài toán nhị phân: precision đo ``đoán dương đúng bao nhiêu''; recall (TPR) đo ``bắt được bao nhiêu dương thật''; specificity (TNR) đo ``bắt được bao nhiêu âm thật''; prevalence = tỉ lệ lớp dương trong tập.}
\kyhieubatdau
  \kh{\mathrm{TP}}{dương thật, đoán dương}
  \kh{\mathrm{TN}}{âm thật, đoán âm}
  \kh{\mathrm{FP}}{âm thật, đoán dương (false alarm)}
  \kh{\mathrm{FN}}{dương thật, đoán âm (bỏ sót)}
  \kh{\mathrm{TPR}}{true positive rate = recall}
  \kh{\mathrm{TNR}}{true negative rate = specificity}
  \kh{\textit{prevalence}}{tỉ lệ mẫu thuộc lớp dương thực tế}
\kyhieuketthuc

\paragraph{Hồi quy (Bài tập 08).}
Phần dư:
\[
e_i=y_i-\hat y_i.
\]
\ynghia{Sai số dự báo tại mẫu \(i\); dùng làm cơ sở cho các metric hồi quy (quy ước đề có thể đảo dấu).}
\kyhieubatdau
  \kh{y_i}{giá trị thực}
  \kh{\hat y_i}{giá trị dự báo}
  \kh{e_i}{phần dư (residual)}
\kyhieuketthuc

\[
SSE=\sum_i e_i^2,\quad
MAE=\frac1n\sum_i|e_i|,\quad
MSE=\frac1n\sum_i e_i^2,\quad
RMSE=\sqrt{MSE}.
\]
\ynghia{Đo độ lớn sai số: SSE/MSE/RMSE nhạy với lỗi lớn (bình phương); MAE trung bình trị tuyệt đối, robust hơn với ngoại lai.}
\kyhieubatdau
  \kh{SSE}{sum of squared errors --- tổng bình phương phần dư}
  \kh{MAE}{mean absolute error}
  \kh{MSE}{mean squared error}
  \kh{RMSE}{căn bậc hai của MSE, cùng đơn vị với \(y\)}
  \kh{n}{số mẫu}
\kyhieuketthuc

RelMSE và CV (tập kiểm tra):
\[
\mathrm{RelMSE}=\frac{\sum_i(y_i-\hat y_i)^2}{\sum_i(y_i-\bar y_{\mathrm{train}})^2},\qquad
CV=\frac{RMSE}{\bar y_{\mathrm{train}}}.
\]
\ynghia{RelMSE so sánh SSE với baseline dự đoán hằng \(\bar y_{\mathrm{train}}\) (càng nhỏ càng tốt, \(<1\) tốt hơn baseline). CV (coefficient of variation) chuẩn hoá RMSE theo mức trung bình của \(y\).}
\kyhieubatdau
  \kh{\bar y_{\mathrm{train}}}{trung bình \(y\) trên tập huấn luyện (mẫu số RelMSE theo đề)}
  \kh{RMSE}{root mean squared error trên tập đang đánh giá}
\kyhieuketthuc

\[
R^2=1-\frac{SSE}{SST},\quad SST=\sum_i(y_i-\bar y)^2.
\]
\ynghia{Hệ số xác định: phần trăm biến thiên của \(y\) được mô hình giải thích so với baseline trung bình; \(R^2=1\) hoàn hảo, \(R^2=0\) ngang baseline.}
\kyhieubatdau
  \kh{R^2}{coefficient of determination}
  \kh{SSE}{tổng bình phương phần dư của mô hình}
  \kh{SST}{total sum of squares --- tổng bình phương lệch so với \(\bar y\)}
  \kh{\bar y}{trung bình \(y\) trên tập đang tính}
\kyhieuketthuc
